\section{Introducción}

La resonancia magnética funcional (fMRI) permite estudiar cambios asociados a la actividad cerebral mediante la señal BOLD. Esta señal no mide directamente la descarga neuronal, sino que refleja una respuesta vascular y metabólica con dinámica temporal propia. Por ello, el análisis de fMRI requiere un modelo que relacione el estímulo experimental con la respuesta hemodinámica observada.

El enfoque más utilizado es el modelo lineal general, en el cual los eventos se convolucionan con una HRF canónica y luego se ajustan a la señal observada. Esta representación ha sido fundamental para el análisis de fMRI, pero supone una forma temporal común que no necesariamente captura variaciones locales o individuales \citep{Friston1998,Glover1999}. Además, se han descrito no linealidades y efectos de saturación que pueden limitar la aproximación de convolución lineal en determinados paradigmas \citep{Friston2000}.

Una alternativa consiste en representar explícitamente la dinámica neurovascular mediante el modelo Balloon--Windkessel, que describe la evolución del flujo sanguíneo, volumen venoso y contenido de desoxihemoglobina \citep{Buxton1998,Friston2000}. Sin embargo, estimar sus estados y parámetros a partir de una señal BOLD ruidosa corresponde a un problema inverso potencialmente mal condicionado.

Las redes neuronales informadas por la física incorporan ecuaciones diferenciales en la función de pérdida, combinando observaciones con conocimiento mecanístico \citep{Raissi2019}. En este proyecto se estudia si esta estrategia permite estimar una HRF y parámetros efectivos con mayor robustez que un GLM y una MLP puramente basada en datos.

\subsection{Pregunta de investigación}

¿Puede una PINN que incorpore el modelo Balloon--Windkessel estimar la respuesta hemodinámica y predecir bloques BOLD no utilizados durante el entrenamiento con mayor robustez e interpretabilidad que un GLM canónico y una MLP sin restricciones fisiológicas?

\subsection{Objetivo general}

Desarrollar y evaluar una PINN inversa para estimar la HRF y parámetros efectivos del modelo Balloon--Windkessel a partir de señales fMRI BOLD, comparando su desempeño con un GLM y una MLP en datos sintéticos y reales.

\subsection{Objetivos específicos}

\begin{enumerate}
    \item Implementar el modelo Balloon--Windkessel y generar señales sintéticas con parámetros conocidos y ruido estructurado.
    \item Diseñar una PINN que combine error de datos, residuos fisiológicos y condiciones iniciales.
    \item Comparar GLM, MLP y PINN mediante evaluación fuera de muestra y recuperación de la HRF.
    \item Aplicar el método a datos motores del Human Connectome Project.
    \item Analizar la consistencia de la HRF y de los parámetros inferidos, además de sus limitaciones.
\end{enumerate}

\subsection{Contribuciones del trabajo}

La contribución principal es una metodología de validación en dos niveles. Primero, se comprueba la recuperabilidad de la HRF y de parámetros conocidos en un entorno sintético. Segundo, se evalúa generalización real mediante entrenamiento y prueba en bloques distintos. Esta separación evita confundir un buen ajuste dentro de muestra con identificación fisiológica o capacidad predictiva.
